\documentclass[10pt,a4paper]{article}
\usepackage[utf8]{inputenc}
\usepackage[T1]{fontenc}
\usepackage[margin=1.8cm]{geometry}
\usepackage{enumitem}
\usepackage{booktabs}
\usepackage{parskip}
\setlength{\parskip}{0.35em}

\title{\textbf{AI-Powered Smart Agriculture Advisory System for Smallholder Farmers in Uganda}\\[0.3em]
\large Response to Research Projects I Poster Feedback}
\author{Project Team}
\date{}

\begin{document}
\maketitle

\textbf{Scope note:} This report critically engages with technical feedback received during the poster session and documents how the project was defended, refined, and improved.

\section*{Selected Technical Comments and Responses}

\textbf{Comment 1 (Faculty Reviewer):} \emph{``The disease model appears to force predictions even for unsupported crop leaves.''}

\textbf{Context:} This directly affects the core objective of accurate image-based disease diagnosis and the reliability of the prototype.

\textbf{Response (Accepted):} The critique was valid. Early inference behavior could still produce overconfident predictions on out-of-distribution inputs, which risked misleading users.

\textbf{Changes made:}
\begin{itemize}[noitemsep]
  \item Added strict upload validation (file type, size, readability, minimum resolution).
  \item Added unsupported-input handling: no diagnosis is returned when confidence is insufficient.
  \item Added crop-selection gate and crop-family consistency checks in the API.
  \item Added crop-conditioned fallback for borderline realistic inputs to reduce unnecessary rejection.
\end{itemize}

\textbf{Comment 2 (Second-Year Reviewer):} \emph{``The weather module sometimes misses key fields like wind speed, so farmers cannot fully trust the advisory.''}

\textbf{Context:} This applies to the weather integration objective and the practical usability of advisory outputs.

\textbf{Response (Accepted):} Valid and relevant. Missing weather attributes reduced trust and weakened the practical value of recommendations.

\textbf{Changes made:}
\begin{itemize}[noitemsep]
  \item Standardized weather response schema to always include wind speed and units.
  \item Integrated OpenWeatherMap (primary) and Open-Meteo fallback for real data continuity.
  \item Updated frontend weather card to correctly render wind speed with API-provided units.
\end{itemize}

\textbf{Comment 3 (Industry Guest Reviewer):} \emph{``Market data looked synthetic; justify realism or use real references.''}

\textbf{Context:} This applies to the market-information objective and credibility of recommendations.

\textbf{Response (Accepted):} Valid. Purely synthetic market values were insufficient for a realistic advisory workflow.

\textbf{Changes made:}
\begin{itemize}[noitemsep]
  \item Replaced random mock generation with realistic regional reference configuration.
  \item Implemented deterministic, reproducible history generation from configured baselines.
  \item Preserved API contracts so frontend charts continued working without regressions.
\end{itemize}

\textbf{Comment 4 (Methodology Panel Reviewer):} \emph{``Your preprocessing notebook mixes disease, market, and RAG, making the disease data pipeline unclear.''}

\textbf{Context:} This concerns methodological clarity, reproducibility, and alignment between documented workflow and actual model training.

\textbf{Response (Accepted):} Valid. The notebook required tighter scope to maintain methodological clarity and reproducibility.

\textbf{Changes made:}
\begin{itemize}[noitemsep]
  \item Refactored notebook to disease-image preprocessing only.
  \item Added defensive checks so cells run cleanly even when optional files/packages are unavailable.
  \item Removed market/RAG executable cells from this notebook and replaced with explicit scope notes.
\end{itemize}

\textbf{Comment 5 (Technical Chair):} \emph{``Your rejection thresholds became too strict after safety updates; real-world photos may be rejected excessively.''}

\textbf{Context:} This affects both model robustness and deployment feasibility for smallholder users with varying camera quality.

\textbf{Response (Partially Accepted):} The concern was valid, but complete removal of safety checks was not acceptable due to misdiagnosis risk. A calibrated middle ground was adopted.

\textbf{Changes made:}
\begin{itemize}[noitemsep]
  \item Recalibrated confidence/margin/entropy thresholds to reduce over-rejection.
  \item Added a cautious-result path (with warning) for plausible but low-confidence cases.
  \item Kept hard rejection for clearly unsupported or unreliable inputs.
\end{itemize}

\section*{Critical Reflection}
The poster feedback significantly improved the project in three ways. First, it increased \textbf{rigour}: uncertain predictions are now handled explicitly, and preprocessing documentation is aligned to the actual training pipeline. Second, it improved \textbf{feasibility}: the system now balances safety and usability through calibrated inference and fallback logic rather than binary accept/reject decisions. Third, it improved \textbf{clarity}: architecture, data scope, and module behavior are now easier to explain and defend during evaluation.

Overall, the feedback shifted the project from a feature-complete prototype toward a more defensible and deployment-aware system for Ugandan smallholder farming contexts.

\end{document}

